\documentclass{article}
\usepackage[spanish]{babel}
\usepackage[pdftex]{graphicx,color} 
\usepackage[utf8]{inputenc}
\usepackage{listings}
\usepackage{colortbl}
\usepackage{textcomp}
\usepackage{listings}
\usepackage{color}
\usepackage{wrapfig}
\usepackage{amsfonts}
\usepackage{amssymb}
\usepackage{algorithmic}
\usepackage{algorithm}
\usepackage[export]{adjustbox}
%\usepackage{amslatex}
\definecolor{lbcolor}{rgb}{0.9,0.9,0.9}
\lstset{
        backgroundcolor=, %\color{lbcolor},
        tabsize=4,
        rulecolor=,
        language=C,
        basicstyle=\scriptsize,
        upquote=true,
        aboveskip={1.5\baselineskip},
        %columns=fixed,
        showstringspaces=false,
        extendedchars=true,
        breaklines=true,
        prebreak = \raisebox{0ex}[0ex][0ex]{\ensuremath{\hookleftarrow}},
        frame=single,
        showtabs=false,
        showspaces=false,
        showstringspaces=false,
        identifierstyle=\ttfamily,
        keywordstyle=\color[rgb]{0,0,1},
        commentstyle=\color[rgb]{0.133,0.545,0.133},
        stringstyle=\color[rgb]{0.627,0.126,0.941},
	literate={á}{{\'a}}1
	{ã}{{\~a}}1
	{é}{{\'e}}1
	{ó}{{\'o}}1
	{í}{{\'i}}1
	{ñ}{{\~n}}1
	{¡}{{!`}}1
	{¿}{{?`}}1
	{ú}{{\'u}}1
	{Í}{{\'I}}1
	{Ó}{{\'O}}1
}

\topmargin -0.5in
\textwidth 6.8in
\textheight 9.0in
\oddsidemargin 0.0in
\evensidemargin 0.0in

\usepackage{enumitem}
\usepackage{multirow}
\usepackage{hhline}
\usepackage{multicol}
%ointpoints{punto}{puntos}

\begin{document}
\vspace{.5pc}
\centerline{\sc Clase Activa 2}
\centerline{\sc ELO320 - Estructura de Datos y Algoritmos}
\centerline{\sc Prof. Nicolás Gálvez Ramírez}
\vspace{2pc}
%\noindent {\bf Instrucciones}:
%\begin{itemize}
%    \item Conteste cada pregunta en una hoja separada. Éstas serán entregadas de tal forma al finalizar el certamen.
%    \item Si decide no contestar una pregunta, debe entregar una hoja sin respuestas en su lugar. 
%    \item Rotule claramente cada una de sus hojas con su Nombre Completo, Rol USM y Número de Paralelo.
%    \item El certamen tiene una duración de 60 minutos.
%\end{itemize}

%\noindent {\bf Preguntas}
%\medskip
\begin{itemize}

\begin{minipage}[h]{0.45\textwidth}
\item \textbf{Definición de Función: Implementación}:
\begin{lstlisting}[language=C, basicstyle=\scriptsize, escapeinside={|}{|}]
float potencia(float a, float b) { 
	return power(a,b);
}
// se define antes de usar


int main(){
	//...
	calculo = x * potencia(y, 2.5);
	//...
}
\end{lstlisting}
\end{minipage}
\hspace{0.5cm}
\begin{minipage}[h]{0.45\textwidth}
\item \textbf{Definición de Función: Prototipo + Implementación}:
\begin{lstlisting}[language=C, basicstyle=\scriptsize, escapeinside={|}{|}]
float potencia(float a, float b);
// prototipo se define antes de usar


int main(){
	//...
	calculo = x * potencia(y, 2.5);
	//...
}

float potencia(float a, float b) { 
	return power(a,b);
}
//implementación puede ir en cualquier parte
//incluso en otro archivo
\end{lstlisting}
\end{minipage}

\begin{minipage}[h]{0.45\textwidth}
\item \textbf{Macro}: No es una variable, es instruccion de preprocesamiento.
\begin{lstlisting}[language=C, basicstyle=\scriptsize, escapeinside={|}{|}]
#include<stdio.h>
#define MACRO 10 

int main(){
	int a = 5;
	printf("|\%|"s\n", a*MACRO);
}
\end{lstlisting}
\end{minipage}
\hspace{0.5cm}
\begin{minipage}[h]{0.45\textwidth}
\item \textbf{Variable Constante}: Inmutable, no se puede modificar.
\begin{lstlisting}[language=C, basicstyle=\scriptsize, escapeinside={|}{|}]

#include<stdio.h>

int main(){
	const int a = 5;
	a = 10; //¿Cómo actúa el compilador?
	printf("|\%|d\n", a);
}
\end{lstlisting}
\end{minipage}

\begin{minipage}[h]{0.45\textwidth}
\item \textbf{Ámbito}
Segmento en el cuál son creadas, existen y pueden ser usadas las variables $\rightarrow$ entre \{ y \}.

\begin{lstlisting}[language=C, basicstyle=\scriptsize, escapeinside={|}{|}] 
#include<stdio.h>

void saludo(){ // var i sirve sólo dentro de función saludo.
	int i = 5;
	printf("Hola alumno numero |\%|d", i);
}

int main(){// var a sirve sólo dentro de función main. 
	int a = 10;
	saludo();
	printf("|\%|d\n", a);
	return 0;
}
\end{lstlisting}
\end{minipage}

\begin{minipage}[h]{0.45\textwidth}
\item \textbf{Variable Local} 
Aquella que está definida dentro de un ámbito determinado, vale decir, entre llaves \{ y \}. Esto aplica para todo elemento que contenga llaves, incluida funciones y estructuras de control. 

\begin{lstlisting}[language=C, basicstyle=\scriptsize, escapeinside={|}{|}]  
#include<stdio.h>

int main(){// var a sirve sólo dentro de función main. 
	int a = 10;

	if(a==10){
		int b = 1;
		printf("|\%|d |\%|d\n", a, b);
	}
	// ¿Qué pasa con la variable b?
	printf("|\%|d |\%|d\n", a, b); 

	return 0;
}
\end{lstlisting}
\end{minipage}
\hspace{0.5cm}
\begin{minipage}[h]{0.45\textwidth}
\item \textbf{Variable Global}
Una variable global, es aquella que está definida \textbf{fuera de cualquier} ámbito. Puede ser llamada, modificada y utilizada desde cualquier ámbito, siempre y cuando \textbf{no exista una variable local} con el mismo nombre. 
\begin{lstlisting}[language=C, basicstyle=\scriptsize, escapeinside={|}{|}] 
#include<stdio.h>

int a=1; //global
int b=2; //global

int main(){
	int a = 10; //local
	b++;
	printf("|\%|d |\%|d\n", a, b);//¿Qué imprime? 
	return 0;
}
\end{lstlisting}
\end{minipage}

\end{itemize}
\end{document}

